\chapter{第十九章：判断力——人类最后的护城河}

# 第十九章：判断力——人类最后的护城河

\textit{"判断力是智力的最高表现形式。"}
——伊曼努尔·康德

\hline\newpage
---

老王我有个朋友，老王我叫他老王，老王是个老股民的。2024年，他跟我说了一件事，差点没把我笑死。

他说他儿子刚毕业进了投行，天天用AI分析股票，说AI比人准。他爸说"儿子啊，爸炒股三十年，什么没见过"，他儿子说"爸你那叫经验主义，现在都AI时代了"。老王一气之下，跟他儿子打赌：各自选五只股票，放一年看谁收益高。

结果年底一算账，老王选的五只，涨了40%。他儿子选的，亏了15%。

他儿子不服气，说"爸你是不是作弊了"。老王说我作弊什么，我就看新闻联播、看身边发生的事、然后凭感觉判断。他儿子说"你这不科学"。老王说"科学有什么用，你那个AI给你分析了半天，最后亏钱了。我凭感觉，反而赚钱了。你说谁的判断力强？"

这当然是个段子，不能当真。但它说明了一个道理：在信息爆炸的时代，判断力比信息本身更值钱。

\section{一、一个投资公司的实验}

## 一、一个投资公司的实验

某家大型投资公司在2024年做了一个实验：他们让AI和一个经验丰富的人类投资经理，同时对一批投资机会做出判断。

AI的分析速度是压倒性的——它可以在几分钟内分析完1000家公司的财务数据、行业趋势、竞争格局。人类投资经理需要两周才能完成同样的工作。

但实验的结果出人意料：当AI和人类投资经理的判断不一致时，最终被证明是正确的判断，有65%是来自人类投资经理，而不是AI。

这个结果让这家公司的管理层开始重新思考：在投资决策中，什么能力是真正重要的？

康德说"判断力是智力的最高表现形式"——这话我琢磨了很久才懂。什么意思？是你聪明不一定有判断力，但有判断力的一定聪明。判断力是把智慧用到实际决策中的能力，这比单纯脑子好使难多了。

\section{二、判断力是什么}

## 二、判断力是什么

判断力是本书反复出现的一个主题。在AI时代，它变得前所未有地重要，值得专门用一章来讨论。

什么是判断力？

判断力不是计算能力。计算能力是处理明确的输入、产生精确输出的能力。计算机在这方面的能力早已超越人类。

判断力不是经验积累。经验是过去经历的沉淀。经验让人类在熟悉的情况下可以快速反应，但它不能帮助人类在全新的情况下做出好的决策。

判断力是\textbf{在不确定性中，决定什么是最重要、以及什么行动方向最可能达到目标的能力}。

它包括几个核心要素：

\begin{itemize}
\item \textbf{情境感知}：理解当前情况的特点和约束
\item \textbf{价值判断}：在多个可能的目标中做出取舍
\item \textbf{因果推理}：理解行动和结果之间的关系
\item \textbf{不确定性处理}：在信息不完整的情况下做出决定

\end{itemize}
就像你站在十字路口，你知道你要去一个地方，但不知道哪条路最近、最不堵、最安全。这不是计算问题——GPS能算出来。但如果你没有GPS、也没有地图、只有一张不知道过期没有的旧地图，你得靠判断：你得看天色、看车流、问路人，然后决定走哪条路。

\section{三、为什么判断力在AI时代变得更稀缺}

## 三、为什么判断力在AI时代变得更稀缺

判断力从来都是稀缺的。但AI让它变得更稀缺了。

原因在于：当AI承担了越来越多的分析性工作时，分析能力本身变得不那么稀缺了。而判断力作为分析能力的"上游"，其稀缺性反而上升了。

更具体地说：

\textbf{当AI可以提供越来越准确的分析时，人们就越需要在分析的基础上做出好的判断。}

如果你问AI"这家公司过去五年的营收增长率是多少"，AI会给你一个准确的答案。但如果你问AI"这家公司是否值得投资"，AI给你的只是一个建议——而你需要判断这个建议是否合理、是否应该采纳。

这个"判断AI的建议是否合理"的能力，就是判断力。当AI越来越强时，对这个能力的需求也越来越高。

就像你会用导航了，但"什么时候该听导航的、什么时候该按自己的判断走"，这件事还是得你自己判断。导航说往前走，但前面堵死了；导航说绕路，但你不熟悉那条路。信导航还是信自己？这就是判断力。

\section{四、判断力的四个来源}

## 四、判断力的四个来源

判断力不是凭空产生的。它有四个主要来源。

\textbf{来源一：经验}

经验是判断力的基础。当你见过足够多的类似情况时，你可以快速识别当前情况与过去的相似之处，并调用过去的经验来指导判断。

但经验也有其局限：它只在熟悉的情况下有效。当遇到全新的情况时，经验可能误导判断。

老王我有个朋友，每次相亲都要问他妈的意见。为什么？因为他妈相过很多次亲，经验丰富。结果有一次，他妈给的意见是"这姑娘不错"，他结了婚三个月就离了——他妈的经验在新时代不适用了。

\textbf{来源二：知识}

知识为判断力提供了背景理解。当你对一个领域有足够的知识时，你可以更好地理解信息的相关性和重要性。

但知识本身不是判断力。拥有大量知识的人，仍然可能做出糟糕的判断——因为他们缺乏把知识转化为判断的能力。

就像你读了一千本游泳书，但你第一次下水还是得喝几口水。知识告诉你怎么游，但判断告诉你什么时候该换气、什么时候该加速、什么时候该悠着点。

\textbf{来源三：思维框架}

思维框架是帮助你在复杂情况下组织思考的结构。拥有好的思维框架，可以帮助你在面对复杂问题时找到正确的分析角度。

不同的思维框架适用于不同类型的问题。选择正确的思维框架，本身就是一种判断。

\textbf{来源四：价值观}

当判断涉及价值取舍时，价值观就成为判断的基准。在多个同样合理但相互冲突的目标之间做选择，需要价值判断。

这种判断没有客观标准——它反映的是一个人或一个组织真正在乎的是什么。

\section{五、判断力的培养}

## 五、判断力的培养

判断力是可以被培养的。以下是几个经过验证的方法。

\textbf{方法一：积累多元经验}

判断力需要经验作为素材。但不是所有的经验都有价值——只有经过反思的经验才有价值。

有意识地积累多元经验，并在每次经历后进行反思：我做对了什么？做错了什么？下次如何改进？这种持续的反思，可以让经验转化为判断力。

\textbf{方法二：刻意练习思维框架}

掌握多种思维框架，并在不同的情况下练习调用正确的框架。

当你面对一个新的问题时，不要急于下结论，先问自己：这是一个什么问题？应该用什么框架来分析？

这种方法需要刻意练习，但随着时间的推移，判断的速度和质量都会提升。

\textbf{方法三：培养认知谦逊}

认知谦逊是承认自己的判断可能是错误的意愿。

具备认知谦逊的人，更容易接受反馈，更容易从错误中学习。而那些过度自信的人，往往会坚持自己最初的判断，即使证据表明他们可能是错的。

《论语》讲"知之为知之，不知为不知，是知也"——知道就是知道，不知道就是不知道，这才叫智慧。这话说的就是认知谦逊：你知道自己不知道什么，这本身就是一种判断力。

\section{六、AI时代的判断力挑战}

## 六、AI时代的判断力挑战

AI时代给人类的判断力带来了新的挑战。

\textbf{挑战一：AI输出的"自信"}

AI系统通常以非常自信的方式输出它们的建议。即使是错误的输出，AI也不会表现出不确定性。这种"自信"可能误导人类用户，让他们过度信任AI的建议。

就像你跟一个特别自信的人聊天，他说得头头是道，你不由自主就信了。但他可能是错的——只是他说话的方式让你觉得他是专家。AI也是这样，它不说"我不确定"，它只说"根据我的分析，答案是A"。

\textbf{挑战二：判断力的"用进废退"}

当AI承担越来越多的决策任务时，人类做出判断的机会在减少。这可能导致判断力的"用进废退"——当真正需要人类判断时，人类已经丧失了这种能力。

就像你天天开车用导航，突然有一天手机没电了，你发现你连回家的路都不记得了——因为你太依赖导航了，你自己判断路径的能力退化了。

\textbf{挑战三：人机协作中的判断责任}

当AI参与决策过程时，判断的责任变得复杂。人类应该接受AI的建议吗？应该在什么情况下override AI的建议？这些问题没有简单的答案，需要人类有清晰的判断力来应对。

\section{七、结语：判断力是AI时代最重要的能力}

## 七、结语：判断力是AI时代最重要的能力

在AI时代，技能经济学的重构在加速。越来越多的技能正在被AI替代。但判断力是少数几个AI难以替代的人类能力之一。

原因在于：判断力需要在不确定性中做出决定，而AI不喜欢不确定性。AI可以处理概率，但它不会像人类那样"感到"不确定性带来的张力。

那些拥有强大判断力的人，将成为AI时代最有价值的成员。他们能够在AI提供的信息和建议的基础上，做出更好的决定。

判断力，是人类在AI时代最后的护城河之一。

\hline\newpage
---

\textbf{本章小结：}

1. 判断力是在不确定性中决定什么是最重要、以及什么行动方向最可能达到目标的能力
2. 判断力的四个来源：经验（基础）、知识（背景）、思维框架（结构）、价值观（取舍标准）
3. AI让判断力变得更稀缺：当分析工作被AI承担后，判断AI建议是否合理的能力需求上升
4. 判断力培养方法：积累多元经验、刻意练习思维框架，培养认知谦逊
5. AI时代的判断力挑战：AI输出的"自信"误导人类、判断力"用进废退"、人机协作中判断责任模糊
6. 判断力是AI时代最重要的能力之一——它是人类最后的护城河之一
